\section{第~\ref{sec:muscles}~章　通往肌肉的输出}

输出的结构——运动单位、突触的安全系数、按大小启动、成对的肌肉、牵拉环路和节律发生器——都已直接测量；有延迟环路的稳定条件是一条定理；身体的内部模型是有充分依据的假说；图中的数字是按本书模型计算的。

{\footnotesize
\setlength{\tabcolsep}{3pt}\renewcommand{\arraystretch}{1.15}
\begin{longtable}{>{\raggedright}p{0.5cm}>{\raggedright}p{4.0cm}>{\raggedright}p{5.6cm}>{\raggedright}p{2.3cm}>{\raggedright\arraybackslash}p{1.7cm}}
\toprule
序号 & 命题 & 实验及测量内容 & 文献 & 状态 \\
\midrule\endfirsthead
\toprule
序号 & 命题 & 实验及测量内容 & 文献 & 状态 \\
\midrule\endhead
1 & 运动单位：一个\nt{ACET}神经元控制一组纤维，从几百根到两千来根；在精细调节的肌肉中单位较小 & 人的肌肉，计数运动轴突和肌纤维：每个神经元平均约有$340$根纤维（手的骨间肌）、约$410$根（肱桡肌）、约$1930$根（腓肠肌内侧头）。本章没有给出最小单位的确切数字。 & \cite{Feinstein1955mu} & 已测量 \\
2 & 肌肉上的突触释放的递质比使纤维兴奋所需的多好几倍 & 神经肌肉突触测量的综述：成年哺乳动物的安全系数（释放的递质与阈值所需之比）通常为$3$--$5$；在人身上，安全系数更多依靠接收方膜的褶皱，而不是释放量。 & \cite{WoodSlater2001} & 已测量 \\
3 & 抽搐~\eqref{eq:twitch}在量级为$50\ms$的时间$T_c$内上升 & 人手的单个运动单位，自主用力，按该单位的脉冲对力作平均：上升时间$30$--$100\ms$，$80\,\%$的单位小于$70\ms$。函数$(t/T_c)e^{1-t/T_c}$是运动单位池模型中的近似。 & \cite{MilnerBrown1973rec, Fuglevand1993} & 已测量 \\
4 & 抽搐叠加~\eqref{eq:forcesum}：$5$、$15$和$40$次/秒时，力分别为$0.2$--$1.1F_1$、$1.8$--$2.2F_1$和约$5.4F_1$（图~\ref{fig:tetanus}） & 按\texttt{models/\allowbreak muscle\_\allowbreak models.py}计算，$T_c=50\ms$：最后$0.3\,\text{s}$内为$0.21$--$1.10$、$1.76$--$2.19$和$5.28$--$5.49\,F_1$。线性叠加是简化：模型中没有真实肌肉的饱和。 & \texttt{models/\allowbreak muscle\_\allowbreak models.py} & 模型 \\
5 & 单位按大小启动；最弱的比最强的弱一百倍 & 猫的腹根：反射输入逐渐增强时，在腹根中脉冲较小的神经元，即小神经元，最先放电。人手骨间肌：各单位的抽搐力从$0.1$到$10$\,g，并随该单位启动时的力近似线性增长。 & \cite{Henneman1957, MilnerBrown1973rec} & 已测量 \\
6 & 力的步长与力成正比，$\Delta F/F\approx q-1$~\eqref{eq:sizeprinciple}：浮点 & 本章由力的等比数列推出。与实验相符：用力大时，同样的力增量所需启动的新单位要少得多。 & 本章推导；\cite{MilnerBrown1973rec} & 理论 \\
7 & 力的离散程度与力本身成正比 & 拇指肌肉的自主等长用力：力的标准差随平均力线性增长；同样的用力若由电刺激引起，离散程度不变。运动单位池模型：按力的顺序启动单位导致线性增长。 & \cite{Jones2002sdn} & 已测量 \\
8 & 动作的平滑性是信号依赖噪声的结果 & 模型：在这种噪声下使终点离散最小的轨迹，再现了实测的眼动和手臂轨迹。尚无把这一原因与其他原因区分开来的直接实验。 & \cite{HarrisWolpert1998} & 假说 \\
9 & 一对肌肉的力之差决定力矩，力之和决定刚度~\eqref{eq:antagonists} & 通过共同收缩进行阻抗控制的理论。人的手臂：用小推力测得的每个关节的刚度大致与该关节的力矩成正比；不同的共同收缩比例改变手的刚度椭圆的形状和方向。 & \cite{Hogan1984, GomiOsu1998stiff} & 已测量 \\
10 & 牵拉传感器兴奋自己的\nt{ACET}神经元，并经抑制性中间神经元关闭拮抗肌（图~\ref{fig:antagonists}） & 猫的脊髓：由牵拉传感器纤维兴奋的单个中间神经元，在拮抗肌的\nt{ACET}神经元中产生单突触抑制电位，突触延迟$0.28$--$0.42\ms$。该实验中没有确定中间神经元的递质（甘氨酸）。 & \cite{Jankowska1972Ia} & 已测量 \\
11 & 环路延迟：脊髓$25$--$30\ms$，经皮层长一倍半到三倍，经眼睛$100$--$200\ms$ & 人的手臂，对关节施加推力：肌肉的短潜伏期响应在$20$--$45\ms$，长潜伏期响应在$45$--$105\ms$，自主响应在$120$--$180\ms$。运动中移动手指的可见位置：平均$160\ms$后作出校正。 & \cite{Pruszynski2008rapid, SaundersKnill2003vis} & 已测量 \\
12 & 当$Gd<\pi/2$时，$\dd x/\dd t=-Gx(t-d)$稳定~\eqref{eq:delaystable}；在边界上频率为$1/(4d)$ & 关于时滞方程根的定理。用\texttt{models/\allowbreak muscle\_\allowbreak models.py}检验，$d=30\ms$：到$3\,\text{s}$时，$G=40$、$50$、$55$每秒对应的幅度为$3\cdot10^{-6}$、$0.13$、$38$（边界为$52$）。 & \cite{Hayes1950} & 理论 \\
13 & 生理性震颤$8$--$12$\,Hz；牵拉环路是其来源之一 & 测量综述：健康人都有约$10$\,Hz范围的细微震颤；其来源是混合的——中枢振荡、单位放电特性、机械共振和反射环路共振，不同条件下主导因素不同。 & \cite{McAuleyMarsden2000} & 假说 \\
14 & 大脑依据身体的内部模型预测指令的结果 & 人用手指捏住重物移动手臂：在惯性、黏性和弹性负载下，握力与负载同时变化，而不是滞后一个环路延迟，也就是说负载被预测到了。综述汇总了这类实验；尚未在神经元中直接观察到模型本身。 & \cite{FlanaganWing1997grip, WolpertGhahramani2000} & 假说 \\
15 & 步行、呼吸、咀嚼的节律由局部网络在恒定输入下产生 & 实验综述：离体神经系统（脊髓、神经节）在没有节律性输入、也没有肌肉反馈的情况下输出节律性的运动模式。 & \cite{MarderBucher2001} & 已测量 \\
16 & 带疲劳的相互抑制~\eqref{eq:halfcentre}产生周期为$0.84\,\text{s}\approx2\tau_a$的节律（图~\ref{fig:halfcentre}） & 定理：在恒定输入下没有稳定状态、带有相互抑制和适应的网络必然振荡。按\texttt{models/\allowbreak muscle\_\allowbreak models.py}计算（$I=1$，$\beta=2$，$b=2$，$\tau=20\ms$，$\tau_a=0.4\,\text{s}$）：周期$0.84\,\text{s}$。真实的发生器正是如此构成，尚未得到证明。 & \cite{Matsuoka1985osc} & 模型 \\
\bottomrule
\end{longtable}}
